\begin{figure}[htbp]
\centering
\begin{tikzpicture}[
    >=Latex,
    font=\small,
    box/.style={
        draw,
        semithick,
        rounded corners=2pt,
        align=center,
        minimum height=10mm,
        minimum width=39mm,
        inner sep=4pt,
        fill=black!2
    },
    key/.style={
        box,
        fill=black!7
    },
    marker/.style={
        draw,
        semithick,
        rounded corners=2pt,
        align=center,
        minimum height=9mm,
        minimum width=22mm,
        inner sep=3pt,
        fill=black!2
    },
    arr/.style={->, semithick},
    note/.style={
        align=center,
        font=\scriptsize,
        text=black!70
    }
]

% ------------------------------------------------------------------
% Titles
% ------------------------------------------------------------------

\node[font=\bfseries\normalsize] at (-4.1,3.45)
    {Locus-wise analysis};

\node[font=\bfseries\normalsize] at (4.1,3.45)
    {Multivariate modeling};

% ------------------------------------------------------------------
% LEFT: LOCUS-WISE
% ------------------------------------------------------------------

\node[key] (lm) at (-4.1,2.45) {
    Genotype matrix\\
    $N$ individuals $\times$ $P$ SNP markers
};

\node[note] at (-4.1,1.80) {
    each SNP is considered separately
};

\node[marker] (s1) at (-6.1,0.85) {SNP$_1$};
\node[marker] (s2) at (-4.1,0.85) {SNP$_2$};
\node[marker] (sp) at (-2.1,0.85) {SNP$_P$};

\node[box] (lst) at (-4.1,-0.70) {
    Marker-specific statistic\\
    e.g.\ $F_{ST}$ or association score
};

\node[key] (lr) at (-4.1,-2.20) {
    SNP ranking\\
    $s_1,\ldots,s_P$
};

\draw[arr] (lm) -- (s1);
\draw[arr] (lm) -- (s2);
\draw[arr] (lm) -- (sp);

\draw[arr] (s1) -- (lst);
\draw[arr] (s2) -- (lst);
\draw[arr] (sp) -- (lst);

\draw[arr] (lst) -- (lr);

% ------------------------------------------------------------------
% RIGHT: MULTIVARIATE
% ------------------------------------------------------------------

\node[key] (rv) at (4.1,2.45) {
    Individual genotype profile\\
    $\mathbf{x}\in\mathbb{R}^{P}$\\
    \scriptsize $P$ SNP values
};

\node[note] at (4.1,1.60) {
    all SNPs are processed jointly
};

\node[box] (rm) at (4.1,0.60) {
    Multivariate / neural model\\
    $f(\mathbf{x})$
};

\node[box] (rz) at (4.1,-0.85) {
    Learned representation\\
    $\mathbf{z}\in\mathbb{R}^{d}$,\quad $d \ll P$
};

\node[key] (ro) at (4.1,-2.35) {
    Breed classification\\
    and/or SNP attribution
};

\draw[arr] (rv) -- (rm);
\draw[arr] (rm) -- (rz);
\draw[arr] (rz) -- (ro);

% ------------------------------------------------------------------
% Separator
% ------------------------------------------------------------------

\draw[gray, dashed] (0,3.10) -- (0,-2.95);

% ------------------------------------------------------------------
% Bottom notation
% ------------------------------------------------------------------

\node[note] at (0,-3.35) {
    $N$: number of individuals \qquad
    $P$: number of SNP markers \qquad
    $d$: latent-space dimension
};

\end{tikzpicture}

\caption{
Conceptual comparison between locus-wise genomic analysis and
multivariate modeling.
Locus-wise approaches evaluate each of the $P$ SNP markers
individually and derive marker-specific statistics, whereas
multivariate models process the complete genotype profile jointly
and may learn a lower-dimensional representation with $d \ll P$.
$N$ denotes the number of individuals in the dataset.
Author's elaboration based on
\cite{weir1984estimating,dimauro2013canonical,bengio2013representation}.
}
\label{fig:inductive_bias_comparison}
\end{figure}